% https://tex.stackexchange.com/a/751662/322482
\documentclass[border=5pt]{standalone}
\usepackage{l3draw}

\ExplSyntaxOn
\cs_new:Npn \intuition_pentagon_point:n #1 {
    \draw_point_polar:nn { 2cm } { 90 + #1 * 360 / 5 }
    % 绘制极坐标上的点 ({90 + #1 * 360 / 5}:2cm)
}
\cs_new:Npn \intuition_pentagon_line:nn #1#2 {
    % 移动到点 #1 绘制线段到点 #2
    \draw_path_moveto:n { \intuition_pentagon_point:n {#1} }
    \draw_path_lineto:n { \intuition_pentagon_point:n {#2} }
}
\ExplSyntaxOff

\begin{document}
\ExplSyntaxOn
\draw_begin:
\draw_scope_begin:
    \draw_set_linewidth:n { .2cm }
    \color_stroke:n { red!25 }
    \int_step_inline:nn { 4 } { % 绘制四条红色的线段
        \intuition_pentagon_line:nn { 0 } { #1 }
        \draw_path_use_clear:n { stroke }
    }
    \intuition_pentagon_line:nn { 2 } { 3 } % 绘制对边的红色线段
    \draw_path_use_clear:n { stroke }
    
    \color_stroke:n { yellow!25 } % 绘制黄色线段
    \intuition_pentagon_line:nn { 1 } { 4 }
    \draw_path_use_clear:n { stroke }
    
    \color_stroke:n { red!50!yellow!25 } % 绘制橙色线段
    \intuition_pentagon_line:nn { 1 } { 2 }
    \draw_path_use_clear:n { stroke }
    
    \color_stroke:n { red!50!blue!25 } % 绘制紫色线段
    \intuition_pentagon_line:nn { 3 } { 4 }
    \draw_path_use_clear:n { stroke }
    
    \color_stroke:n { blue!25 } % 绘制两条蓝色线段
    \intuition_pentagon_line:nn { 1 } { 3 }
    \intuition_pentagon_line:nn { 2 } { 4 }
    \draw_path_use_clear:n { stroke }
\draw_scope_end:

\draw_path_moveto:n { % 回到起始点
    \intuition_pentagon_point:n { 0 }
}
\int_step_inline:nn { 5 } {
    \draw_path_lineto:n {
        \intuition_pentagon_point:n {#1}
    }
}
\int_step_inline:nn { 5 } {
    \draw_path_lineto:n {
        \intuition_pentagon_point:n { #1 * 2 } % 利用了360度的周期性
    }
}
\draw_path_use_clear:n { stroke }

\int_step_inline:nn { 5 } {
    \draw_path_circle:nn { % 绘制黑色圆点
        \intuition_pentagon_point:n {#1}
    } { 6pt }
    \draw_path_use_clear:n { fill }
}
\draw_end:
\ExplSyntaxOff
\end{document}